% §5.3 Structural foundation for the renter-share moderator hypothesis.
% Builds on Kurlat & Stroebel (2015) RFS heterogeneously-informed-buyer
% model + a Gaussian public-signal extension from Vives (2008) Ch. 2.
% This subsection converts the renter-share moderator from rhetorical
% primary-outcome designation to structurally derived prediction.

\subsection{Structural foundation for the renter-share moderator}
\label{sec:structural-renter}

The pooled meta-regression in Section~\ref{sec:meta-regression-12}
identifies the renter share of housing stock as the primary moderator
of the Shen-Ross uniqueness premium under the housing-search
information-frictions literature.  Here we make the directional
prediction structurally precise by extending
\citet{kurlat2015testing}'s heterogeneously-informed-buyer model with
a Gaussian public signal in the \citet{vives2008information} tradition.
Two auxiliary assumptions are required and we state them explicitly.

\paragraph{Setup.} Following \citet{kurlat2015testing}, the valuation
of house $h$ in neighborhood $j$ decomposes as
\begin{equation}
v_h = \beta_h \theta_j + (1 - \beta_h) \eta_h,
\label{eq:ks-decomp}
\end{equation}
where $\theta_j$ is the unobserved quality of neighborhood $j$ and
$\eta_h$ is house-specific idiosyncratic quality.
\citeauthor{kurlat2015testing} distinguish two buyer types: informed
buyers observe a binary signal about $v_h$, while uninformed buyers
do not.  Their Proposition~4 establishes that
$\partial^2 v / (\partial \theta \partial \beta_h) = 1 > 0$, i.e.,
a house's response to a neighborhood-quality shock is increasing in
its loading $\beta_h$.

\paragraph{Auxiliary assumption 1: listing text as Gaussian public
signal of $\theta_j$.}
We treat the listing text content $T_h$ as a noisy public signal of
the neighborhood-quality component:
\begin{equation}
T_h = \theta_j + \xi_h, \qquad \xi_h \sim N(0, \sigma_T^2),
\label{eq:public-signal}
\end{equation}
observed by both buyer types prior to bidding. This places the model
in the normal-linear signal-extraction tradition of
\citet{vives2008information} Chapter 2.  Under
\eqref{eq:public-signal} and a Gaussian prior on $\theta_j$, an
informed buyer with private-signal precision $\tau_I$ on $\theta_j$
forms the posterior
\begin{equation}
E_I[\theta_j \mid T_h, \text{private}] = \alpha_I \, \text{private}
+ (1 - \alpha_I) T_h,
\quad
\alpha_I = \frac{\tau_I}{\tau_I + \tau_T},
\label{eq:posterior-informed}
\end{equation}
where $\tau_T = 1/\sigma_T^2$ is the precision of the public signal.
An uninformed buyer has $\tau_U = 0$, so $\alpha_U = 0$ and her
posterior is simply $T_h$.

\paragraph{Auxiliary assumption 2: renter share as proxy for the
uninformed-buyer share.}
Let $r$ denote the share of uninformed buyers in market $j$.
Following \citet{chinco2016misinformed}, who establish that out-of-
town and recent-mover buyers carry less neighborhood-specific
information than long-term resident buyers, and
\citet{kurlat2015testing}'s own use of recent-mover share as a
proxy for the informed-seller share, we operationalize $r$ as the
metropolitan-area renter share of housing stock. This is the
empirical link in the derivation and we report sensitivity to
alternative buyer-information proxies (recent-mover share, length-
of-residence percentiles) in Appendix~\ref{app:moderator-robustness}.

\paragraph{Equilibrium price.} Under the
\citet{kurlat2015testing} equilibrium pricing rule (their equations
10-11), the marginal-buyer transaction price for observationally
indistinguishable houses is the share-weighted average of the two
buyer types' expected valuations.  Substituting equations
\eqref{eq:ks-decomp}, \eqref{eq:public-signal}, and
\eqref{eq:posterior-informed},
\begin{align}
\log p_h^* &\approx r \cdot E_U[v_h \mid T_h] + (1 - r) \cdot
E_I[v_h \mid T_h] \\
&= r \beta_h T_h + (1 - r) \beta_h
[\alpha_I \, \text{private} + (1 - \alpha_I) T_h]
+ \text{(constants in $\eta_h$)}.
\label{eq:price}
\end{align}

\paragraph{Moderation prediction.} Differentiating
\eqref{eq:price} twice yields
\begin{equation}
\frac{\partial \log p_h^*}{\partial T_h}
= \beta_h \cdot [1 - (1 - r) \alpha_I],
\qquad
\frac{\partial^2 \log p_h^*}{\partial T_h \partial r}
= \beta_h \alpha_I > 0.
\label{eq:moderation}
\end{equation}
The partial effect of listing text on sale price is increasing in the
share of uninformed buyers.  Operationalizing the latter as the
metropolitan-area renter share gives the empirical prediction
\begin{equation}
\frac{\partial^2 \log p_h^*}{\partial T_h \partial \text{(renter share)}_j} > 0,
\label{eq:empirical-prediction}
\end{equation}
which is the primary moderator hypothesis tested in
Section~\ref{sec:meta-regression-12}.

\paragraph{Robustness to the auxiliary assumptions.} The two
auxiliary assumptions are not innocuous and we note their role.
First, the Gaussian public-signal assumption \eqref{eq:public-signal}
imposes a normal-linear signal-extraction structure that
\citeauthor{kurlat2015testing} themselves do not adopt; their buyer
signal is binary.  We follow \citet{vives2008information} in
adopting the Gaussian extension because (i) the binary case is a
limiting case of the Gaussian case as signal precision tends to a
threshold, and (ii) the empirical content of the moderation
prediction in \eqref{eq:moderation} is robust to the specific
distributional form, requiring only that uninformed buyers place
greater weight on the public signal than informed buyers (a property
true of all reasonable signal structures).  Second, the renter-share-
as-uninformed-buyer-share map is empirical; we report sensitivity to
alternative proxies in Appendix~\ref{app:moderator-robustness} and
flag this as a limitation of any attempt to identify buyer-
information moderation from market-level demographic aggregates.
The prediction in \eqref{eq:empirical-prediction} survives under any
proxy that is increasing in the uninformed-buyer share.
